\chapter{Types of Risk}
\label{chap:Types of Risk}

%---------------------------- SECTION ----------------------------
\section{Making Sense of the Risk Landscape}

\paragraph{By this point, we have established what risk is, what a risk assessment is designed to achieve, and how we arrived at the regulatory and professional environment that makes this work so important today. Now it is time to look at the \textit{substance} of risk itself — the different forms it takes, the ways it presents in organisational life, and why understanding the full landscape matters before you begin any assessment process.}

\paragraph{One of the most common shortcomings in organisational risk management is a tendency towards \textbf{tunnel vision} — a focus on the risks that are most familiar, most visible, or most recently in the regulatory spotlight, at the expense of a broader and more complete picture. A financial services firm that focuses exclusively on credit and market risk may overlook operational or reputational exposures. A technology business laser-focused on cybersecurity risk may not give adequate attention to its regulatory compliance obligations or its supply chain vulnerabilities.}

\paragraph{This chapter is designed to help you avoid that trap. It offers a broad taxonomy of risk types — not as a definitive or exhaustive classification system, but as a practical map of the terrain. Think of it as a checklist of \textit{categories to consider} when approaching any risk assessment, rather than a rigid framework to be applied mechanically}

\paragraph{It is also worth noting at the outset that the boundaries between risk categories are rarely clean. In practice, risks interact, overlap, and amplify each other. A cybersecurity incident can become a data protection failure, which becomes a regulatory breach, which becomes a reputational crisis, which becomes a financial loss. Understanding risk types individually is valuable — but understanding how they connect is what separates a good risk assessment from a great one}

%---------------------------- SECTION ----------------------------
\section{A Note on Risk Taxonomy}
\paragraph{A \textbf{risk taxonomy} is simply an organised, consistently defined classification of the types of risk an organisation faces. Having a shared taxonomy matters for the same reason that having a shared language matters — it ensures that people across the organisation are talking about the same things, that risks are captured consistently in registers and reports, and that nothing falls through the gaps because it did not fit neatly into someone's mental category}

\paragraph{Different organisations, industries, and frameworks use different taxonomies, and there is no single universally correct version. What matters is that your organisation's taxonomy is:}
\begin{itemize}
  \item{\textbf{Comprehensive enough} to capture all types of material risk relevant to your business}
  \item{\textbf{Clear enough} to be used consistently by different people and teams}
  \item{\textbf{Practical enough} to actually inform decision-making rather than satisfy documentation requirements}
\end{itemize}

\paragraph{With that in mind, what follows is a broad taxonomy suitable as a starting point for most organisations, with particular relevance for those operating in regulated sectors. You may find that some categories need to be expanded, merged, or adapted to fit your specific context — and that is entirely appropriate.}

%---------------------------- SECTION ----------------------------
\section{Operational Risk}
\paragraph{\textbf{Operational risk} refers to the risk of loss or harm resulting from inadequate or failed internal processes, people, systems, or external events. It is, in many respects, the broadest and most pervasive category of risk — present in virtually every activity an organisation undertakes.}
\paragraph{Examples of operational risk include:}
\begin{itemize}
  \bitem{Human error}{mistakes made by staff in carrying out their duties.}
  \bitem{Process failures}{breakdowns in internal procedures or controls.}
  \bitem{System failures}{technology outages, software errors, or infrastructure failures.}
  \bitem{Fraud and misconduct}{deliberate wrongdoing by employees or third parties.}
  \bitem{Third-party failures}{suppliers, outsourced service providers, or partners failing to perform as expected.}
  \bitem{Physical events}{fire, flood, or other incidents affecting premises or operations.}
\end{itemize}
\paragraph{For compliance professionals, operational risk is particularly significant because many regulatory requirements are, at their core, designed to manage specific types of operational risk. Requirements around controls, record-keeping, staff competence, and outsourcing oversight are all responses to identified operational risk exposures.}
\paragraph{The \textbf{Basel II framework}, introduced in 2004, formally recognised operational risk as a distinct risk category requiring capital allocation in banking — a recognition that has since influenced risk management thinking well beyond the financial services sector.}

%---------------------------- SECTION ----------------------------
\section{Legal and Regulatory Risk}
\paragraph{\textbf{Legal risk} is the risk of financial loss, disruption, reputational damage, or other adverse consequences arising from legal proceedings, the unenforceability of contracts, or changes in the legal environment. \textbf{Regulatory risk} — sometimes treated as a subset of legal risk, sometimes as a separate category — refers specifically to the risk of failing to comply with applicable laws, regulations, rules, and standards, and the consequences that follow.}
\paragraph{Legal risk is the risk of financial loss, disruption, reputational damage, or other adverse consequences arising from legal proceedings, the unenforceability of contracts, or changes in the legal environment. Regulatory risk — sometimes treated as a subset of legal risk, sometimes as a separate category — refers specifically to the risk of failing to comply with applicable laws, regulations, rules, and standards, and the consequences that follow:}
\begin{itemize}
  \item{Entering contracts that are unenforceable or expose the organisation to unexpected liability.}
  \item{Failure to comply with sector-specific regulatory requirements — licencing, business rules, reporting obligations, etc.}
  \item{Changes in legislation or regulation that affect the legality or viability of existing business practices.}
  \item{Regulatory investigations, enforcement actions, fines, or sanctions.}
  \item{Litigation, whether initiated by customers, counterparties, regulators, or other parties.}
  \item{Failure to maintain required documentation or records to evidence compliance.}
\end{itemize}
\paragraph{It is worth noting that regulatory risk has a particularly important characteristic: \textbf{it is not static}. The regulatory environment changes continuously, sometimes gradually and sometimes suddenly, and organisations need risk assessment processes that are capable of identifying and responding to those changes in a timely way. We will return to this theme in Chapter~\ref{chap:Legal and Regulatory Considerations}.}

%---------------------------- SECTION ----------------------------
\section{Financial Risk}
\paragraph{\textbf{Financial risk} encompasses the range of risks that directly affect an organisation's financial position — its income, assets, liabilities, cash flow, and overall financial stability. In its broadest sense, it includes:}
\begin{itemize}
  \bitem{Credit risk}{the risk that a counterparty will fail to meet a financial obligation, such as a borrower defaulting on a loan or a customer failing to pay.}
  \bitem{Market risk}{the risk of losses arising from movements in market prices, interest rates, exchange rates, or commodity prices.}
  \bitem{Liquidity risk}{the risk that an organisation will be unable to meet its financial obligations as they fall due, even if it is technically solvent.}
  \bitem{Capital risk}{the risk that an organisation does not hold sufficient capital to absorb losses or to satisfy regulatory capital requirements.}
\end{itemize}
\paragraph{For organisations in financial services, these categories are subject to detailed regulatory requirements and sophisticated quantitative management techniques. For organisations outside financial services, the concepts remain relevant — a business that does not manage its credit exposure to customers, or that holds insufficient liquidity to meet short-term obligations, faces real and serious consequences regardless of its sector.}

%---------------------------- SECTION ----------------------------
\section{Strategic Risk}
\paragraph{\textbf{Strategic risk} refers to the risk that the organisation's overall strategy — its chosen direction, business model, competitive positioning, or major decisions — will fail to deliver the intended outcomes, or will expose it to unexpected harm.}
\paragraph{Strategic risks are often among the most significant and the hardest to manage, in part because they tend to be driven by factors that are difficult to control — market disruption, shifts in customer behaviour, technological change, competitive dynamics, or macroeconomic developments.}
\paragraph{Examples include:}
\begin{itemize}
  \item{Entering a new market or launching a product that proves unviable.}
  \item{Failing to adapt to a major shift in technology or customer expectation.}
  \item{Making an acquisition that does not deliver anticipated value — or that introduces unexpected liabilities.}
  \item{Pursuing a growth strategy that outpaces the organisation's operational or governance capabilities.}
  \item{A fundamental misalignment between the organisation's stated strategy and its actual risk appetite.}
\end{itemize}
\paragraph{Strategic risk is sometimes underrepresented in compliance-focused risk assessments, which tend to concentrate on more immediate and measurable exposures. However, for organisations where compliance is integral to the business model — regulated financial institutions, healthcare providers, legal services firms — strategic and regulatory risk are deeply intertwined, and strategic decisions can rarely be made without careful consideration of their compliance implications.}

%---------------------------- SECTION ----------------------------
\section{Reputational Risk}
\paragraph{Reputational risk is the risk that an event, action, or association will damage the organisation's standing in the eyes of its customers, clients, regulators, investors, employees, or the wider public — with consequential impacts on its ability to operate effectively and sustainably.}
\paragraph{Reputational risk deserves particular attention for several reasons:}
\begin{itemize}
  \item{It is often the \textbf{consequence} of other risks materialising, rather than a standalone event. A data breach, a regulatory sanction, a product failure, or a governance scandal all carry reputational consequences that can outlast and outweigh their immediate financial or legal impact.}
  \item{It is \textbf{highly contextual} — what damages one organisation's reputation may have little effect on another, depending on their stakeholder base, sector, and public profile.}
  \item{It is \textbf{difficult to quantify} using the standard tools of risk assessment, which can lead to it being underweighted in formal processes despite its practical significance.}
  \item{It can be \textbf{disproportionately rapid} in its development. In an environment where information travels instantly and public scrutiny is intense, reputational damage can escalate from a minor incident to a major crisis in a matter of hours.}
\end{itemize}
\paragraph{For compliance and legal professionals, reputational risk is particularly salient because regulatory breaches and legal failures are among the most common triggers of serious reputational damage. We will dedicate a full chapter to this topic in Chapter~\ref{chap:Reputational Risk Assessment}.}

%---------------------------- SECTION ----------------------------
\section{Cybersecurity and Information Risk}
\paragraph{\textbf{Cybersecurity risk} refers to the risk of harm arising from the failure, compromise, or misuse of digital systems, technology infrastructure, or data. \textbf{Information risk} is the broader category, encompassing risks to data and information regardless of whether they arise from cyber threats specifically.}
\paragraph{\textbf{:}}
\begin{itemize}
  \bitem{Cyber attacks}{ransomware, phishing, distributed denial of service (DDoS), and other malicious activities.}
  \bitem{Data breaches}{unauthorised access to, or disclosure of, personal or sensitive data.}
  \bitem{System failures}{Technology and infrastructure outages affecting business continuity.}
  \bitem{Insider threats}{deliberate or accidental misuse of data or systems by employees or contractors.}
  \bitem{Third-party technology risk}{vulnerabilities introduced through suppliers, software providers, or cloud services.}
\end{itemize}
\paragraph{The regulatory dimension of cybersecurity risk has grown significantly in recent years. Data protection legislation — including the GDPR and its equivalents in other jurisdictions — imposes specific obligations around the security of personal data and requires organisations to assess and manage the risks associated with their data processing activities. Sector-specific regulators in financial services, healthcare, and critical infrastructure have introduced increasingly detailed cybersecurity requirements.}
\paragraph{We will explore cybersecurity risk assessment in depth in Chapter~\ref{chap:Cybersecurity and Data Protection Risk Assessment}}

%---------------------------- SECTION ----------------------------
\section{People and Human Capital Risk}
\paragraph{People risk refers to the risk of harm or loss arising from issues related to the organisation's workforce — including recruitment, retention, conduct, competence, culture, and wellbeing.}
\paragraph{In compliance-sensitive environments, people risk takes on particular significance because the behaviour of individuals — whether deliberate or inadvertent — is one of the most common sources of regulatory and legal exposure. Key examples include:}
\begin{itemize}
  \bitem{Misconduct}{deliberate wrongdoing by employees, including fraud, bribery, market manipulation, or breach of client confidentiality.}
  \bitem{Errors and omissions}{unintentional mistakes that lead to regulatory breaches or client harm.}
  \bitem{Competence gaps}{insufficient knowledge or skill to perform regulated activities appropriately.}
  \bitem{Cultural failures}{an organisational culture that tolerates or incentivises inappropriate behaviour}
  \bitem{Key person dependency}{over-reliance on specific individuals whose departure would create significant operational or knowledge risk.}
  \bitem{Employment law}{claims arising from discrimination, unfair dismissal, or failure to meet employment obligations.}
\end{itemize}
\paragraph{The increasing focus of regulators on \textbf{culture and conduct} — particularly in financial services, but increasingly across other sectors — reflects a growing recognition that people risk is not merely an HR concern. It is a core compliance and governance issue.}

%---------------------------- SECTION ----------------------------
\section{Third-Party and Supply Chain Risk}
\paragraph{\textbf{Third-party risk} refers to the risk that arises from an organisation's relationships with external parties — suppliers, outsourced service providers, contractors, agents, distributors, and other partners. \textbf{Supply chain risk} is the subset of third-party risk concerned with the chain of suppliers and subcontractors involved in delivering goods or services}
\paragraph{Third-party risk refers to the risk that arises from an organisation's relationships with external parties — suppliers, outsourced service providers, contractors, agents, distributors, and other partners. Supply chain risk is the subset of third-party risk concerned with the chain of suppliers and subcontractors involved in delivering goods or services.}
\begin{itemize}
  \bitem{Supplier failure}{a critical supplier ceasing to operate or failing to deliver, disrupting the organisation's own operations.}
  \bitem{3rd Party conduct and compliance failure}{ a supplier engaged in bribery, modern slavery, environmental violations, or data protection breaches creates legal, regulatory, and reputational exposure for the organisations that work with them.}
  \bitem{Outsourcing risk}{where regulated activities are outsourced, the organisation typically retains regulatory responsibility for those activities regardless of who performs them.}
  \bitem{Concentration risk}{over-reliance on a small number of critical suppliers or a single geographic region.}
  \bitem{Data security in 3rd party relationships}{sharing data with third parties creates information security and data protection exposure.}
\end{itemize}
\paragraph{Regulators across multiple sectors — particularly financial services — have significantly tightened expectations around third-party risk management in recent years, reflecting the systemic risks that can arise when critical functions are concentrated in a small number of providers.}

%---------------------------- SECTION ----------------------------
\section{Environmental and Climate Risk}
\paragraph{\textbf{Environmental risk} refers to the risk of harm — to the organisation, to third parties, or to the natural environment — arising from the organisation's impact on, or exposure to, environmental conditions. \textbf{Climate risk} has emerged as a critical subcategory, encompassing both:}
\begin{itemize}
  \bitem{Physical climate risk}{the risk of direct harm from climate-related events such as flooding, extreme weather, or temperature changes affecting operations, assets, or supply chains.}
  \bitem{Transition risk}{the risk arising from the shift towards a lower-carbon economy, including regulatory change, shifts in market demand, technology change, and reputational exposure for organisations perceived as inadequately addressing their environmental impact.}
\end{itemize}
\paragraph{Environmental and climate risk assessment is an area of rapidly growing regulatory focus. Disclosure requirements, mandatory climate-related risk reporting, and the integration of climate risk into prudential regulation are all developing quickly, and organisations in virtually every sector are likely to face increasing expectations in this area. We will explore this in depth in Chapter~\ref{chap:Environmental and Sustainability Risk Assessment}.}

%---------------------------- SECTION ----------------------------
\section{Project and Change Risk}
\paragraph{\textbf{Project risk} refers to the risk that a specific initiative — whether a technology implementation, a business transformation, a product launch, a merger, or a regulatory change programme — will not deliver its intended outcomes, or will create unintended harm in the process.}
\paragraph{\textbf{Change risk} is closely related — the risk that organisational change of any kind, including restructuring, process redesign, or system replacement, will disrupt existing controls, create new vulnerabilities, or introduce compliance gaps.}
\paragraph{For compliance professionals, change risk is particularly significant because major organisational changes frequently have implications for regulatory compliance — and those implications are not always fully considered during the design and implementation phase. A technology migration that inadvertently disrupts a key control, or a restructuring that creates ambiguity about regulatory ownership, can create serious compliance exposure that would have been entirely avoidable with earlier and more thorough risk assessment.}

%---------------------------- SECTION ----------------------------
\section{How Risk Types Interact — The Compound Effect}
\paragraph{As noted at the outset of this chapter, the boundaries between these categories are rarely clean in practice. It is worth illustrating this with a brief example}
\paragraph{Consider an organisation that decides to launch a new digital product targeting retail customers. In doing so, it is simultaneously taking on:}
\begin{itemize}
  \bitem{Strategic risk}{will the product succeed in the market?}
  \bitem{Operational risk}{will the underlying processes and technology work reliably?}
  \bitem{Legal and regulatory risk}{does the product meet all applicable consumer protection, financial promotion, and data protection requirements?}
  \bitem{Cybersecurity risk}{does the platform adequately protect customer data?}
  \bitem{Third-party risk}{are the technology providers and outsourced service partners appropriately managed?}
  \bitem{People risk}{do staff understand their obligations in relation to the new product?}
  \bitem{Reputational risk}{what happens to public perception if any of the above goes wrong?}
  \bitem{Financial risk}{what is the exposure if the product underperforms or generates unexpected liabilities}
\end{itemize}
\paragraph{Each of these risk types exists independently, but they also interact and amplify each other. A cybersecurity failure leads to a data breach, which triggers a regulatory investigation, which generates adverse media coverage, which affects customer confidence, which impacts financial performance. The compound effect of multiple risk types materialising together — or one triggering another — is often far more serious than any individual risk considered in isolation.}
\paragraph{This interconnectedness is one of the strongest arguments for taking a holistic approach to risk assessment — one that looks across the full risk landscape rather than examining each category in isolation.}

%---------------------------- SECTION ----------------------------
\newpage
\section{Chapter Summary}
\paragraph{Before moving on, let us anchor the key ideas from this chapter:}

\paragraph{Key takeaways:}

\sloppy
\begin{itemize}
  \bitem{Operational}{Failures in processes, people, systems, or external events.}
  \bitem{Legal \& Regulatory}{Legal proceedings, unenforceability, and compliance failures.}
  \bitem{Financial}{Credit, market, liquidity, and capital exposures.}
  \bitem{Strategic}{Failures of direction, business model, or major decisions.}
  \bitem{Reputational}{Damage to standing with key stakeholders.}
  \bitem{Cybersecurity \& Information}{Digital threats, data breaches, and technology failures.}
  \bitem{People \& Human Capital}{Conduct, competence, culture, and workforce issues.}
  \bitem{Third-Party \& Supply Chain}{Risks arising from external relationships and dependencies.}
  \bitem{Environmental \& Climate}{Physical and transition risks related to environmental conditions.}
  \bitem{Project \& Change}{Risks arising from initiatives, transformation, and organisational change.}
  \bitem{Takeaways}{In this chapter we learned:
  \begin{itemize}
          \item{No single risk taxonomy is universally correct — what matters is that yours is comprehensive, clear, and consistently applied.}
          \item{Risk types rarely exist in isolation — understanding how they interact is as important as understanding each individually.}
          \item{Some risk types, particularly reputational and strategic risk, are easily underweighted in formal assessments despite their practical significance.}
          \item{The regulatory landscape around several risk categories — cyber, climate, third-party — is evolving rapidly and will continue to generate new compliance obligations.}
          \item{A holistic approach to risk assessment, one that looks across the full landscape, is considerably more effective than one focused narrowly on the most familiar categories.}
        \end{itemize}
    }
\end{itemize}
\fussy

\barquote{In Part II, we will move from the foundations of risk thinking to the frameworks and standards that provide structure to the risk assessment process — exploring what the major frameworks offer, how regulatory expectations are shaped, and how to choose an approach that fits your organisation.}